\documentclass[10pt,a4paper]{amsart}
\usepackage[margin=19mm]{geometry}
\usepackage{amsmath,amssymb,mathtools}
\usepackage[scr=rsfs,cal=euler]{mathalfa}
\usepackage{xfrac}
\usepackage[unicode,hidelinks,bookmarks=false]{hyperref}
\newcommand{\ie}{i.e.\ }
\newcommand{\cf}{cf.\ }
\newcommand{\resp}{respectively\ }
\newcommand{\Subsection}{\subsection}
\providecommand{\nobreakdash}{\nobreak\hbox{-}}
\setlength{\emergencystretch}{2em}
\multlinegap0pt
\usepackage{bm}
\usepackage{bbm}
\usepackage{dsfont}
\usepackage{xspace}
\usepackage{cancel}
\usepackage{mathtools}
\usepackage{amsthm}
\makeatletter
\@ifundefined{theorem}{\newtheorem{theorem}{Theorem}}{}
\makeatother
\title[Small complete 3-term progression free sets in cyclic groups]{Small complete 3-term progression free sets in cyclic groups and vector spaces}
\author{Bence Csajb\'ok, Zolt\'an L\'or\'ant Nagy}
\date{Source: arXiv:2606.30186v1 (CC BY 4.0)}
\begin{document}
\maketitle

\noindent\emph{Source and licence.} Small complete 3-term progression free sets in cyclic groups and vector spaces. Version arXiv:2606.30186v1 (CC BY 4.0), distributed under the Creative Commons Attribution 4.0 International licence, as recorded in the arXiv licence metadata for this exact version and shown on the abstract page https://arxiv.org/abs/2606.30186v1. Full rights notice: licenses/arxiv-2606.30186-CC-BY-4.0.txt.\par\smallskip

\noindent\textit{Editorial scope.} Counted unit: the Theorem whose source label is \texttt{None} in the article (source0.tex lines 965--979, source label \texttt{None}), delivered together with the article's own complete original proof of it (source0.tex lines 981--1102, one \texttt{proof} environment of 122 lines). No mathematics is written, repaired or re-derived: statement and proof are the article's own text line for line. Required context retained uncounted: none. Every definition, display and label of those lines is retained unchanged. Omitted from this package and not redistributed: 1--964, 980--980, 1103--1529. Nothing inside a retained excerpt is omitted or altered: every retained body is delivered exactly as the article's lines are written, so no omission receipt of any kind accompanies this package. Citations: the retained excerpts carry no citation command at all, so no bibliography entry of the article is transcribed and no reference is redistributed. Reference adaptation: none was needed and none was made; every reference inside the delivered unit resolves inside it, so no unresolved reference or citation is produced. Printed slips kept literal: no printed word, symbol, hypothesis or number of the counted statement or of its proof was altered. Layout and class: the article's own class is \texttt{article} with packages including geometry, amsmath, amssymb, amsthm, mathtools, enumitem, booktabs, array, longtable, hyperref, cleveref, todonotes; the wrapper is the lane's pinned \texttt{amsart} editorial wrapper at 10pt on A4, which restates the theorem-like environments the retained text needs and adds the article's own macro definitions that the retained text needs (none), with extra wrapper packages (bm, bbm, dsfont, xspace, cancel, mathtools). The delivered numbering is the unit's own. Assets: no retained span contains a figure environment, a graphics-inclusion command, a PDF-page inclusion, a table or a TikZ picture; no image file is copied into this package, the package declares no asset dependency and no image file is redistributed with it. Licence: version arXiv:2606.30186v1 of this article is distributed under the Creative Commons Attribution 4.0 International licence, as recorded in the arXiv licence metadata for this exact version and shown on the abstract page https://arxiv.org/abs/2606.30186v1. A full rights notice is delivered at licenses/arxiv-2606.30186-CC-BY-4.0.txt. Characters: every retained excerpt is pure ASCII, so the delivered engine, XeLaTeX with its default Latin Modern text font, reports no missing character.
	\begin{theorem}
		Let \(p\) be an odd prime and let \(\varepsilon>0\).  Then there is a constant
		\(C_{p,\varepsilon}\) such that, for every \(n\ge1\),
		\[
		a(3\text{-}\mathrm{AP},\mathbb F_p^n)
		\le
		C_{p,\varepsilon}\,n^{1+\varepsilon}p^{n/2}.
		\]
		In particular, for every fixed odd prime \(p\),
		\[
		a(3\text{-}\mathrm{AP},\mathbb F_p^n)
		=
		p^{n/2+o(n)}.
		\]
	\end{theorem}

	\begin{proof}
		Choose \(N_0=N_0(p,\varepsilon)\) large enough so that the estimates below hold
		for every \(n\ge N_0\).  We then prove the bound for \(n\ge N_0\) by strong
		induction on \(n\), and finally enlarge \(C_{p,\varepsilon}\) to cover the
		finitely many dimensions \(n<N_0\).
		
		Let \(n\) be large.  Choose \(k\ge2\) minimal with
		\[
		p^k>\left(\frac nk\right)^2 .
		\]
		By minimality,
		\[
		p^{k-1}\le \left(\frac{n}{k-1}\right)^2,
		\]
		and hence
		\[
		p^{k/2}
		\le
		\sqrt p\,\frac{n}{k-1}.
		\tag{5.1}
		\]
		
		Choose an odd integer \(L\ge1\) such that
		$ n_0=k(2L+1)\le n$
		is as large as possible.  Since \(L\) is odd, the admissible values of
		\(2L+1\) differ by \(4\).  Thus the remainder
		$ r=n-n_0$
		satisfies
		\[
		0\le r<4k .
		\tag{5.2}
		\]
		Moreover,
		\[
		2L-1<2L+1\le \frac nk,
		\]
		so
		\[
		p^k>\left(\frac nk\right)^2>(2L-1)^2 .
		\]
		Therefore the subfield quadratic graph construction applies.  It gives a two-sided
		complete \(3\)-AP-free set $A\subseteq \mathbb F_p^{n_0}$
		of size
		\[
		|A|
		=
		p^{k(L+1)}
		=
		p^{n_0/2}p^{k/2}.
		\]
		
		If \(r=0\), then by \((5.1)\),
		\[
		|A|
		\le
		\sqrt p\,\frac{n}{k-1}p^{n/2}.
		\]
		For all sufficiently large \(n\), this is at most
		\[
		n^{1+\varepsilon}p^{n/2},
		\]
		and the remaining finitely many cases are absorbed into the constant
		\(C_{p,\varepsilon}\).
		
		Assume now that \(r>0\).  Since \(r<n\), the induction hypothesis gives a
		complete \(3\)-AP-free set
		$ B\subseteq \mathbb F_p^r$
		with
		\[
		|B|\le
		C_{p,\varepsilon}\,r^{1+\varepsilon}p^{r/2}.
		\]
		Since \(A\) is endpoint-midpoint complete  and \(B\) is complete \(3\)-AP-free, the
		product lemma gives that
		\[
		A\times B
		\subseteq
		\mathbb F_p^{n_0}\times\mathbb F_p^r
		\cong
		\mathbb F_p^n
		\]
		is complete \(3\)-AP-free.  Its size is at most
		\[
		|A||B|
		\le
		C_{p,\varepsilon}\,
		r^{1+\varepsilon}p^{k/2}p^{n/2}.
		\]
		Using \((5.1)\) and \((5.2)\), we get
		\[
		r^{1+\varepsilon}p^{k/2}
		<
		(4k)^{1+\varepsilon}
		\sqrt p\,\frac{n}{k-1}.
		\]
		For all sufficiently large \(n\), the right-hand side is at most
		$   n^{1+\varepsilon}.$
		Indeed, \(k=O_p(\log n)\), so
		\[
		(4k)^{1+\varepsilon}\sqrt p\,\frac{n}{k-1}
		=
		O_{p,\varepsilon}\!\left(n(\log n)^\varepsilon\right)
		=
		o\!\left(n^{1+\varepsilon}\right).
		\]
		Thus, for all sufficiently large \(n\),
		\[
		|A||B|
		\le
		C_{p,\varepsilon}\,n^{1+\varepsilon}p^{n/2}.
		\]
		This completes the induction for \(n\ge N_0\).  Enlarging
		\(C_{p,\varepsilon}\) to cover the finitely many dimensions \(n<N_0\) proves
		the theorem for all \(n\ge1\).
		
		Finally,
		\[
		n^{1+\varepsilon}p^{n/2}=p^{n/2+o(n)}
		\]
		for fixed \(p\), and since \(\varepsilon>0\) was arbitrary, the last assertion
		follows.
	\end{proof}

\end{document}
